\section{Analyse et conception}
\label{sec:conception}

Notre système se divise en deux parties bien distinctes : d'un côté la chaîne du modèle, qui tourne sur Colab, et de l'autre le simulateur, qui tourne en local. Le simulateur ne réentraîne jamais le modèle : il se contente de charger ses poids et ses métadonnées (figure~\ref{fig:architecture}).

\begin{figure}[H]
  \centering
  % Architecture fonctionnelle globale d'AutoMaze AI (cinq étages, de l'environnement au contrôleur)
\begin{tikzpicture}[
    etage/.style={draw=accent, fill=accent!7, rounded corners=3pt, align=center, font=\footnotesize,
                  text width=12.2cm, minimum height=1.0cm, inner sep=4pt},
    bloc/.style={draw=black!60, fill=white, rounded corners=3pt, align=center, font=\footnotesize,
                 text width=3.6cm, minimum height=1.05cm, inner sep=3pt},
    ml/.style={bloc, draw=accent, fill=accent!18, very thick},
    fleche/.style={-{Stealth[length=2.2mm]}, draw=black!70, thick}
  ]
  \node[etage] (env) at (0,0) {\textbf{1. Environnement de simulation}\\Routes · Véhicule · Objets routiers · Destination};

  \node[bloc] (camera) at (-4.3,-1.75) {\textbf{Caméra virtuelle}\\Image de l'objet};
  \node[ml] (modele) at (-4.3,-3.35) {\textbf{Modèle ML}\\Classification de l'objet};
  \node[bloc] (infos) at (4.3,-2.55) {\textbf{Informations du simulateur}\\Coordonnées et positions des objets};
  \node[bloc] (etat) at (0,-4.75) {\textbf{État des obstacles}\\Position + nature identifiée};

  \node[etage] (fusion) at (0,-6.3) {\textbf{2. Fusion des informations et représentation de l'environnement}\\Association de la classe prédite aux coordonnées de l'objet correspondant};
  \node[etage] (plan) at (0,-7.75) {\textbf{3. Planificateur}\\Graphe · Dijkstra · Recalcul du trajet};
  \node[etage] (decision) at (0,-9.2) {\textbf{4. Système de décision}\\Ralentir · Arrêter · Contourner · Continuer};
  \node[etage] (controle) at (0,-10.65) {\textbf{5. Contrôleur du véhicule}\\Vitesse · Freinage · Direction};
  \node[font=\footnotesize\bfseries, text=integration] (but) at (0,-11.75) {Déplacement autonome jusqu'à la destination};

  \draw[fleche] (env.south -| camera) -- (camera);
  \draw[fleche] (camera) -- (modele);
  \draw[fleche] (env.south -| infos) -- (infos);
  \draw[fleche] (modele.south) |- (etat.west);
  \draw[fleche] (infos.south) |- (etat.east);
  \draw[fleche] (etat) -- (fusion);
  \draw[fleche] (fusion) -- (plan);
  \draw[fleche] (plan) -- (decision);
  \draw[fleche] (decision) -- (controle);
  \draw[fleche] (controle) -- (but);

  % Une décision qui bloque un tronçon relance le planificateur
  \draw[fleche, dashed, draw=accent] (decision.east) -- ++(0.55,0) |- node[pos=0.25, right, font=\scriptsize, align=left, text=accent] {obstacle\\bloquant :\\recalcul} (plan.east);
  % Le véhicule se déplace dans l'environnement : la boucle reprend
  \draw[fleche, draw=integration] (but.west) -- ++(-6.9,0) |- node[pos=0.25, left, font=\scriptsize, text=integration, align=right] {nouvelle\\position} (env.west);

  % Ce que le simulateur ne transmet pas
  \node[font=\scriptsize, text=black!65, align=center, text width=3.8cm] at (4.3,-3.95) {La vraie classe de l'objet n'est jamais transmise à la décision};
\end{tikzpicture}

  \caption{Architecture : le simulateur ne dépend du modèle que par ses poids et ses métadonnées.}
  \label{fig:architecture}
\end{figure}

Dans le simulateur, les trois classifieurs possibles fonctionnent de la même façon : ils reçoivent une image et renvoient une classe, un score et 6 probabilités. Nous pouvons donc passer de l'un à l'autre sans toucher à la navigation.

\textbf{Analyse fonctionnelle.} Concrètement, le modèle reçoit le recadrage d'un objet et renvoie une classe avec un score. Si ce score est trop faible, le simulateur affiche \og Inconnu \fg{} et la voiture s'arrête par prudence. Nous choisirons ce seuil sur le jeu de validation.

\textbf{Acteurs.} L'équipe s'occupe des données, de l'entraînement et de l'évaluation des modèles. L'encadrant valide nos choix au fil du projet, et le jury juge à la fin les résultats et la démonstration.

\subsection*{Règles de décision proposées}

\begin{center}
\begin{tabularx}{\textwidth}{>{\raggedright\arraybackslash}p{2.8cm}>{\raggedright\arraybackslash}p{4.2cm}Y}
  \toprule
  \textbf{Classe prédite} & \textbf{Décision} & \textbf{Effet sur l'itinéraire} \\
  \midrule
  Personne & S'arrêter et attendre & Tronçon bloqué quelques pas, puis la voiture repart \\
  Vélo, moto & Ralentir & Coût du tronçon multiplié \\
  Voiture & Contourner & Tronçon bloqué, nouveau calcul \\
  Bus, camion & Contourner largement & Tronçon bloqué, tronçons voisins pénalisés, nouveau calcul \\
  Inconnu & Arrêt prudent & Tronçon bloqué, nouveau calcul \\
  \bottomrule
\end{tabularx}
\end{center}

Dans le simulateur, chaque obstacle se comporte selon sa vraie classe : une personne finit de traverser, alors qu'un camion reste immobile. Si le modèle se trompe, la voiture peut donc prendre une mauvaise décision, voire provoquer une collision, et nous pouvons compter ces erreurs. Ces règles sont une première proposition, que nous devons encore valider en équipe.
